% --- Archon physics formalization source begin ---
% archon:physics
% archon:covers PhyXMiniProblems/problem_phyx_mini_0849.lean
% archon:source-report reports/phyx_mini/problem_phyx_mini_0849.source.json

\chapter{Physics problem phyx\_mini\_0849}
\label{ch:PhyXMiniProblems_problem_phyx_mini_0849}

\paragraph{Problem source.}
\#\# Physical scenario

The conducting box in figure has been given an excess negative charge. The surface density of excess electrons at the center of the top surface is \textbackslash{}( 5.0 	imes 10\textasciicircum{}\{10\} \textbackslash{}) electrons/m\textbackslash{}(\textasciicircum{}2\textbackslash{}).

\#\# Question

What are the electric field strengths \textbackslash{}( E\_1 \textbackslash{}) to \textbackslash{}( E\_3 \textbackslash{}) at points 1 to 3?

\#\# Answer choices

- A: A: 800N m\textasciicircum{}2/C

- B: B: 1000N m\textasciicircum{}2/C

- C: C: 200N m\textasciicircum{}2/C

- D: D: 900N m\textasciicircum{}2/C

\#\# Figure caption (auxiliary; use the image as primary evidence)

The image consists of two rectangular shapes, one inside the other. 

1. The outer rectangle is light gray and encompasses the entire area.
2. The inner rectangle is white and is centered within the outer rectangle.
3. Three numbered black dots are present:
   - Dot labeled "1" is above both rectangles.
   - Dot labeled "2" is inside the gray area but outside the white rectangle.
   - Dot labeled "3" is inside the white rectangle.

There is a clear distinction between the two areas with the inner rectangle appearing as a highlighted or focused region within the outer one.

\#\# Dataset metadata

Electrostatics; Physical Model Grounding Reasoning, Spatial Relation Reasoning

\paragraph{Recorded answer/context.}
D: D: 900N m\textasciicircum{}2/C

\paragraph{Figure/image path.}
/root/proposal\_for\_physic/science-mango/phyx\_mini\_run/phyx\_data/test\_image/849.png

\paragraph{Formalization target.}
create a compiling Lean file with sorry bodies at `PhyXMiniProblems/problem\_phyx\_mini\_0849.lean`. The Lean declarations must preserve the physical quantities, dimensions or dimensional roles, figure labels, governing-law hypotheses, and final relation expressed by this problem.
Use Mathlib/Physlib names found through LeanExplore where available. If a physics API is missing, introduce faithful local abstractions rather than scalar placeholder aliases.

\begin{theorem}[Physics formalization target]
\label{thm:physics:phyx_mini_0849:target}
\lean{PhyXMiniProblems.ProblemPhyXMini0849.problem_phyx_mini_0849}
\uses{def:physics:phyx-mini-0849:phyxminiproblems-problemphyxmini0849-conductingboxsetup, def:physics:phyx-mini-0849:phyxminiproblems-problemphyxmini0849-matchesconductingboxscenario, def:physics:phyx-mini-0849:phyxminiproblems-problemphyxmini0849-matchesreportedsurfacedensitydata, def:physics:phyx-mini-0849:phyxminiproblems-problemphyxmini0849-matchessuppliedconductingboxfigure, def:physics:phyx-mini-0849:phyxminiproblems-problemphyxmini0849-usesstandardelectrostaticconstants, def:physics:phyx-mini-0849:phyxminiproblems-problemphyxmini0849-hasphysicalconductingboxparameters, def:physics:phyx-mini-0849:phyxminiproblems-problemphyxmini0849-satisfiesconductingboxelectrostatics, def:physics:phyx-mini-0849:phyxminiproblems-problemphyxmini0849-electricfieldstrengthinnewtonspercoulomb, def:physics:phyx-mini-0849:phyxminiproblems-problemphyxmini0849-recordeddatasetanswer, def:physics:phyx-mini-0849:phyxminiproblems-problemphyxmini0849-roundstonearestresolution, def:physics:phyx-mini-0849:phyxminiproblems-problemphyxmini0849-isuniquematchinganswer}
The assigned autoformalize agent should translate this physics problem into Lean declarations in the covered file, with theorem and lemma proof bodies written as `by sorry`.
\end{theorem}
\begin{proof}
This is an autoformalization task, not a proof task. Produce statements that can later be proved by the physics prover without weakening the physical model.
\end{proof}

% --- Archon declaration topology begin ---
\section{Declaration topology}
\begin{definition}[surface Number Density Dimension]
  \label{def:physics:phyx-mini-0849:phyxminiproblems-problemphyxmini0849-surfacenumberdensitydimension}
  \lean{PhyXMiniProblems.ProblemPhyXMini0849.surfaceNumberDensityDimension}
  The dimension L\textasciicircum{}-2 of a surface number density
\end{definition}
\begin{proof}
  This is the declared model definition of surface Number Density Dimension.
\end{proof}
\begin{definition}[surface Charge Density Dimension]
  \label{def:physics:phyx-mini-0849:phyxminiproblems-problemphyxmini0849-surfacechargedensitydimension}
  \lean{PhyXMiniProblems.ProblemPhyXMini0849.surfaceChargeDensityDimension}
  The dimension C L\textasciicircum{}-2 of a signed surface charge density
\end{definition}
\begin{proof}
  This is the declared model definition of surface Charge Density Dimension.
\end{proof}
\begin{definition}[vacuum Permittivity Dimension]
  \label{def:physics:phyx-mini-0849:phyxminiproblems-problemphyxmini0849-vacuumpermittivitydimension}
  \lean{PhyXMiniProblems.ProblemPhyXMini0849.vacuumPermittivityDimension}
  The dimension C\textasciicircum{}2 T\textasciicircum{}2 M\textasciicircum{}-1 L\textasciicircum{}-3 of vacuum permittivity
\end{definition}
\begin{proof}
  This is the declared model definition of vacuum Permittivity Dimension.
\end{proof}
\begin{definition}[electric Field Dimension]
  \label{def:physics:phyx-mini-0849:phyxminiproblems-problemphyxmini0849-electricfielddimension}
  \lean{PhyXMiniProblems.ProblemPhyXMini0849.electricFieldDimension}
  The dimension M L T\textasciicircum{}-2 C\textasciicircum{}-1 of an electric field
\end{definition}
\begin{proof}
  This is the declared model definition of electric Field Dimension.
\end{proof}
\begin{definition}[Surface Electron Number Density Quantity]
  \label{def:physics:phyx-mini-0849:phyxminiproblems-problemphyxmini0849-surfaceelectronnumberdensityquantity}
  \lean{PhyXMiniProblems.ProblemPhyXMini0849.SurfaceElectronNumberDensityQuantity}
  \uses{def:physics:phyx-mini-0849:phyxminiproblems-problemphyxmini0849-surfacenumberdensitydimension}
  A nonnegative physical count of electrons per unit area
\end{definition}
\begin{proof}
  This is the declared model definition of Surface Electron Number Density Quantity.
\end{proof}
\begin{definition}[Charge Magnitude Quantity]
  \label{def:physics:phyx-mini-0849:phyxminiproblems-problemphyxmini0849-chargemagnitudequantity}
  \lean{PhyXMiniProblems.ProblemPhyXMini0849.ChargeMagnitudeQuantity}
  A nonnegative physical magnitude of electric charge
\end{definition}
\begin{proof}
  This is the declared model definition of Charge Magnitude Quantity.
\end{proof}
\begin{definition}[Signed Surface Charge Density Quantity]
  \label{def:physics:phyx-mini-0849:phyxminiproblems-problemphyxmini0849-signedsurfacechargedensityquantity}
  \lean{PhyXMiniProblems.ProblemPhyXMini0849.SignedSurfaceChargeDensityQuantity}
  \uses{def:physics:phyx-mini-0849:phyxminiproblems-problemphyxmini0849-surfacechargedensitydimension}
  A signed physical surface charge density
\end{definition}
\begin{proof}
  This is the declared model definition of Signed Surface Charge Density Quantity.
\end{proof}
\begin{definition}[Vacuum Permittivity Quantity]
  \label{def:physics:phyx-mini-0849:phyxminiproblems-problemphyxmini0849-vacuumpermittivityquantity}
  \lean{PhyXMiniProblems.ProblemPhyXMini0849.VacuumPermittivityQuantity}
  \uses{def:physics:phyx-mini-0849:phyxminiproblems-problemphyxmini0849-vacuumpermittivitydimension}
  A nonnegative physical vacuum permittivity
\end{definition}
\begin{proof}
  This is the declared model definition of Vacuum Permittivity Quantity.
\end{proof}
\begin{definition}[Electric Field Vector Quantity]
  \label{def:physics:phyx-mini-0849:phyxminiproblems-problemphyxmini0849-electricfieldvectorquantity}
  \lean{PhyXMiniProblems.ProblemPhyXMini0849.ElectricFieldVectorQuantity}
  \uses{def:physics:phyx-mini-0849:phyxminiproblems-problemphyxmini0849-electricfielddimension}
  A physical electric-field vector in three-dimensional space
\end{definition}
\begin{proof}
  This is the declared model definition of Electric Field Vector Quantity.
\end{proof}
\begin{definition}[surface Electron Number Density Per Square Meter]
  \label{def:physics:phyx-mini-0849:phyxminiproblems-problemphyxmini0849-surfaceelectronnumberdensitypersquaremeter}
  \lean{PhyXMiniProblems.ProblemPhyXMini0849.surfaceElectronNumberDensityPerSquareMeter}
  \uses{def:physics:phyx-mini-0849:phyxminiproblems-problemphyxmini0849-surfaceelectronnumberdensityquantity}
  Coherent-SI readout in electrons per square metre
\end{definition}
\begin{proof}
  This is the declared model definition of surface Electron Number Density Per Square Meter.
\end{proof}
\begin{definition}[charge Magnitude In Coulombs]
  \label{def:physics:phyx-mini-0849:phyxminiproblems-problemphyxmini0849-chargemagnitudeincoulombs}
  \lean{PhyXMiniProblems.ProblemPhyXMini0849.chargeMagnitudeInCoulombs}
  \uses{def:physics:phyx-mini-0849:phyxminiproblems-problemphyxmini0849-chargemagnitudequantity}
  Coherent-SI readout of a charge magnitude in coulombs
\end{definition}
\begin{proof}
  This is the declared model definition of charge Magnitude In Coulombs.
\end{proof}
\begin{definition}[signed Surface Charge Density In Coulombs Per Square Meter]
  \label{def:physics:phyx-mini-0849:phyxminiproblems-problemphyxmini0849-signedsurfacechargedensityincoulombspersquaremeter}
  \lean{PhyXMiniProblems.ProblemPhyXMini0849.signedSurfaceChargeDensityInCoulombsPerSquareMeter}
  \uses{def:physics:phyx-mini-0849:phyxminiproblems-problemphyxmini0849-signedsurfacechargedensityquantity}
  Coherent-SI readout of signed surface charge in coulombs per square metre
\end{definition}
\begin{proof}
  This is the declared model definition of signed Surface Charge Density In Coulombs Per Square Meter.
\end{proof}
\begin{definition}[vacuum Permittivity In Farads Per Meter]
  \label{def:physics:phyx-mini-0849:phyxminiproblems-problemphyxmini0849-vacuumpermittivityinfaradspermeter}
  \lean{PhyXMiniProblems.ProblemPhyXMini0849.vacuumPermittivityInFaradsPerMeter}
  \uses{def:physics:phyx-mini-0849:phyxminiproblems-problemphyxmini0849-vacuumpermittivityquantity}
  Coherent-SI readout of vacuum permittivity in farads per metre
\end{definition}
\begin{proof}
  This is the declared model definition of vacuum Permittivity In Farads Per Meter.
\end{proof}
\begin{definition}[electric Field Vector In Newtons Per Coulomb]
  \label{def:physics:phyx-mini-0849:phyxminiproblems-problemphyxmini0849-electricfieldvectorinnewtonspercoulomb}
  \lean{PhyXMiniProblems.ProblemPhyXMini0849.electricFieldVectorInNewtonsPerCoulomb}
  \uses{def:physics:phyx-mini-0849:phyxminiproblems-problemphyxmini0849-electricfieldvectorquantity}
  Coherent-SI readout of an electric-field vector in newtons per coulomb
\end{definition}
\begin{proof}
  This is the declared model definition of electric Field Vector In Newtons Per Coulomb.
\end{proof}
\begin{definition}[physlib Elementary Charge Magnitude In Coulombs]
  \label{def:physics:phyx-mini-0849:phyxminiproblems-problemphyxmini0849-physlibelementarychargemagnitudeincoulombs}
  \lean{PhyXMiniProblems.ProblemPhyXMini0849.physlibElementaryChargeMagnitudeInCoulombs}
  Physlib provides the elementary charge as a ChargeUnit, rather than as a dimensionful charge quantity. Its ratio to the coulomb calibrates the independent physical magnitude in the setup
\end{definition}
\begin{proof}
  This is the declared model definition of physlib Elementary Charge Magnitude In Coulombs.
\end{proof}
\begin{definition}[Figure Point]
  \label{def:physics:phyx-mini-0849:phyxminiproblems-problemphyxmini0849-figurepoint}
  \lean{PhyXMiniProblems.ProblemPhyXMini0849.FigurePoint}
  The three numbered black dots in image 849
\end{definition}
\begin{proof}
  This is the declared model definition of Figure Point.
\end{proof}
\begin{definition}[Figure Region]
  \label{def:physics:phyx-mini-0849:phyxminiproblems-problemphyxmini0849-figureregion}
  \lean{PhyXMiniProblems.ProblemPhyXMini0849.FigureRegion}
  The qualitative region occupied by a labelled dot in the image
\end{definition}
\begin{proof}
  This is the declared model definition of Figure Region.
\end{proof}
\begin{definition}[Surface Sample Location]
  \label{def:physics:phyx-mini-0849:phyxminiproblems-problemphyxmini0849-surfacesamplelocation}
  \lean{PhyXMiniProblems.ProblemPhyXMini0849.SurfaceSampleLocation}
  The surface location at which the electron density is reported
\end{definition}
\begin{proof}
  This is the declared model definition of Surface Sample Location.
\end{proof}
\begin{definition}[Excess Charge Sign]
  \label{def:physics:phyx-mini-0849:phyxminiproblems-problemphyxmini0849-excesschargesign}
  \lean{PhyXMiniProblems.ProblemPhyXMini0849.ExcessChargeSign}
  The sign of the box's excess charge
\end{definition}
\begin{proof}
  This is the declared model definition of Excess Charge Sign.
\end{proof}
\begin{definition}[Electrostatic Model]
  \label{def:physics:phyx-mini-0849:phyxminiproblems-problemphyxmini0849-electrostaticmodel}
  \lean{PhyXMiniProblems.ProblemPhyXMini0849.ElectrostaticModel}
  The electrostatic idealization used for the conducting box
\end{definition}
\begin{proof}
  This is the declared model definition of Electrostatic Model.
\end{proof}
\begin{definition}[Conducting Box Figure]
  \label{def:physics:phyx-mini-0849:phyxminiproblems-problemphyxmini0849-conductingboxfigure}
  \lean{PhyXMiniProblems.ProblemPhyXMini0849.ConductingBoxFigure}
  \uses{def:physics:phyx-mini-0849:phyxminiproblems-problemphyxmini0849-figurepoint, def:physics:phyx-mini-0849:phyxminiproblems-problemphyxmini0849-figureregion}
  Literal qualitative information carried by image 849. Gray material is the conductor and the centered white rectangle is its cavity
\end{definition}
\begin{proof}
  This is the declared model definition of Conducting Box Figure.
\end{proof}
\begin{definition}[Conducting Box Setup]
  \label{def:physics:phyx-mini-0849:phyxminiproblems-problemphyxmini0849-conductingboxsetup}
  \lean{PhyXMiniProblems.ProblemPhyXMini0849.ConductingBoxSetup}
  \uses{def:physics:phyx-mini-0849:phyxminiproblems-problemphyxmini0849-surfaceelectronnumberdensityquantity, def:physics:phyx-mini-0849:phyxminiproblems-problemphyxmini0849-chargemagnitudequantity, def:physics:phyx-mini-0849:phyxminiproblems-problemphyxmini0849-signedsurfacechargedensityquantity, def:physics:phyx-mini-0849:phyxminiproblems-problemphyxmini0849-vacuumpermittivityquantity, def:physics:phyx-mini-0849:phyxminiproblems-problemphyxmini0849-electricfieldvectorquantity, def:physics:phyx-mini-0849:phyxminiproblems-problemphyxmini0849-figurepoint, def:physics:phyx-mini-0849:phyxminiproblems-problemphyxmini0849-surfacesamplelocation, def:physics:phyx-mini-0849:phyxminiproblems-problemphyxmini0849-excesschargesign, def:physics:phyx-mini-0849:phyxminiproblems-problemphyxmini0849-electrostaticmodel, def:physics:phyx-mini-0849:phyxminiproblems-problemphyxmini0849-conductingboxfigure}
  Independent physical data for the box. No requested field strength, answer choice, or rounding interval is stored in this setup
\end{definition}
\begin{proof}
  This is the declared model definition of Conducting Box Setup.
\end{proof}
\begin{definition}[Matches Conducting Box Scenario]
  \label{def:physics:phyx-mini-0849:phyxminiproblems-problemphyxmini0849-matchesconductingboxscenario}
  \lean{PhyXMiniProblems.ProblemPhyXMini0849.MatchesConductingBoxScenario}
  \uses{def:physics:phyx-mini-0849:phyxminiproblems-problemphyxmini0849-conductingboxsetup}
  Qualitative physical assumptions stated or implied by the problem
\end{definition}
\begin{proof}
  This is the declared model definition of Matches Conducting Box Scenario.
\end{proof}
\begin{definition}[Matches Reported Surface Density Data]
  \label{def:physics:phyx-mini-0849:phyxminiproblems-problemphyxmini0849-matchesreportedsurfacedensitydata}
  \lean{PhyXMiniProblems.ProblemPhyXMini0849.MatchesReportedSurfaceDensityData}
  \uses{def:physics:phyx-mini-0849:phyxminiproblems-problemphyxmini0849-surfaceelectronnumberdensitypersquaremeter, def:physics:phyx-mini-0849:phyxminiproblems-problemphyxmini0849-conductingboxsetup}
  The sole numerical datum in the question: the top-center density is 5.0 * 10\textasciicircum{}10 excess electrons per square metre
\end{definition}
\begin{proof}
  This is the declared model definition of Matches Reported Surface Density Data.
\end{proof}
\begin{definition}[Matches Supplied Conducting Box Figure]
  \label{def:physics:phyx-mini-0849:phyxminiproblems-problemphyxmini0849-matchessuppliedconductingboxfigure}
  \lean{PhyXMiniProblems.ProblemPhyXMini0849.MatchesSuppliedConductingBoxFigure}
  \uses{def:physics:phyx-mini-0849:phyxminiproblems-problemphyxmini0849-conductingboxsetup}
  Primary-image readout of the outer box, inner cavity, and three dots
\end{definition}
\begin{proof}
  This is the declared model definition of Matches Supplied Conducting Box Figure.
\end{proof}
\begin{definition}[Uses Standard Electrostatic Constants]
  \label{def:physics:phyx-mini-0849:phyxminiproblems-problemphyxmini0849-usesstandardelectrostaticconstants}
  \lean{PhyXMiniProblems.ProblemPhyXMini0849.UsesStandardElectrostaticConstants}
  \uses{def:physics:phyx-mini-0849:phyxminiproblems-problemphyxmini0849-chargemagnitudeincoulombs, def:physics:phyx-mini-0849:phyxminiproblems-problemphyxmini0849-vacuumpermittivityinfaradspermeter, def:physics:phyx-mini-0849:phyxminiproblems-problemphyxmini0849-physlibelementarychargemagnitudeincoulombs, def:physics:phyx-mini-0849:phyxminiproblems-problemphyxmini0849-conductingboxsetup}
  Calibration of the independent dimensionful constants. The elementary charge is tied to Physlib's unit-scale declaration. Vacuum permittivity is tied to Physlib's positive FreeSpace.ε₀ and its standard SI value. None of these fields contains a requested electric-field answer
\end{definition}
\begin{proof}
  This is the declared model definition of Uses Standard Electrostatic Constants.
\end{proof}
\begin{definition}[Has Physical Conducting Box Parameters]
  \label{def:physics:phyx-mini-0849:phyxminiproblems-problemphyxmini0849-hasphysicalconductingboxparameters}
  \lean{PhyXMiniProblems.ProblemPhyXMini0849.HasPhysicalConductingBoxParameters}
  \uses{def:physics:phyx-mini-0849:phyxminiproblems-problemphyxmini0849-surfaceelectronnumberdensitypersquaremeter, def:physics:phyx-mini-0849:phyxminiproblems-problemphyxmini0849-chargemagnitudeincoulombs, def:physics:phyx-mini-0849:phyxminiproblems-problemphyxmini0849-vacuumpermittivityinfaradspermeter, def:physics:phyx-mini-0849:phyxminiproblems-problemphyxmini0849-conductingboxsetup}
  Nondegeneracy conditions needed by the boundary calculation
\end{definition}
\begin{proof}
  This is the declared model definition of Has Physical Conducting Box Parameters.
\end{proof}
\begin{definition}[Satisfies Conducting Box Electrostatics]
  \label{def:physics:phyx-mini-0849:phyxminiproblems-problemphyxmini0849-satisfiesconductingboxelectrostatics}
  \lean{PhyXMiniProblems.ProblemPhyXMini0849.SatisfiesConductingBoxElectrostatics}
  \uses{def:physics:phyx-mini-0849:phyxminiproblems-problemphyxmini0849-surfaceelectronnumberdensitypersquaremeter, def:physics:phyx-mini-0849:phyxminiproblems-problemphyxmini0849-chargemagnitudeincoulombs, def:physics:phyx-mini-0849:phyxminiproblems-problemphyxmini0849-signedsurfacechargedensityincoulombspersquaremeter, def:physics:phyx-mini-0849:phyxminiproblems-problemphyxmini0849-vacuumpermittivityinfaradspermeter, def:physics:phyx-mini-0849:phyxminiproblems-problemphyxmini0849-electricfieldvectorinnewtonspercoulomb, def:physics:phyx-mini-0849:phyxminiproblems-problemphyxmini0849-conductingboxsetup}
  The governing electrostatic laws used by the solution: excess electrons give the signed density sigma = -n e; immediately outside an equilibrated conductor, E = (sigma/epsilon\_0) n; the field vanishes in conducting material at equilibrium; an enclosed charge-free cavity is electrostatically shielded. The laws are generic over points having the relevant region role. They do not state the requested numerical field strengths at points 1--3
\end{definition}
\begin{proof}
  This is the declared model definition of Satisfies Conducting Box Electrostatics.
\end{proof}
\begin{definition}[electric Field Strength In Newtons Per Coulomb]
  \label{def:physics:phyx-mini-0849:phyxminiproblems-problemphyxmini0849-electricfieldstrengthinnewtonspercoulomb}
  \lean{PhyXMiniProblems.ProblemPhyXMini0849.electricFieldStrengthInNewtonsPerCoulomb}
  \uses{def:physics:phyx-mini-0849:phyxminiproblems-problemphyxmini0849-electricfieldvectorinnewtonspercoulomb, def:physics:phyx-mini-0849:phyxminiproblems-problemphyxmini0849-figurepoint, def:physics:phyx-mini-0849:phyxminiproblems-problemphyxmini0849-conductingboxsetup}
  Field strength at a labelled dot, in newtons per coulomb
\end{definition}
\begin{proof}
  This is the declared model definition of electric Field Strength In Newtons Per Coulomb.
\end{proof}
\begin{definition}[Answer Choice]
  \label{def:physics:phyx-mini-0849:phyxminiproblems-problemphyxmini0849-answerchoice}
  \lean{PhyXMiniProblems.ProblemPhyXMini0849.AnswerChoice}
  The four answer labels printed in the source
\end{definition}
\begin{proof}
  This is the declared model definition of Answer Choice.
\end{proof}
\begin{definition}[displayed Exterior Field Strength In Newtons Per Coulomb]
  \label{def:physics:phyx-mini-0849:phyxminiproblems-problemphyxmini0849-displayedexteriorfieldstrengthinnewtonspercoulomb}
  \lean{PhyXMiniProblems.ProblemPhyXMini0849.displayedExteriorFieldStrengthInNewtonsPerCoulomb}
  \uses{def:physics:phyx-mini-0849:phyxminiproblems-problemphyxmini0849-answerchoice}
  The answer strings contain the dimensionally inconsistent text N m\textasciicircum{}2/C. Because the question asks for electric-field strength, their numeric values are modeled with the physically compatible unit N/C
\end{definition}
\begin{proof}
  This is the declared model definition of displayed Exterior Field Strength In Newtons Per Coulomb.
\end{proof}
\begin{definition}[displayed Resolution In Newtons Per Coulomb]
  \label{def:physics:phyx-mini-0849:phyxminiproblems-problemphyxmini0849-displayedresolutioninnewtonspercoulomb}
  \lean{PhyXMiniProblems.ProblemPhyXMini0849.displayedResolutionInNewtonsPerCoulomb}
  \uses{def:physics:phyx-mini-0849:phyxminiproblems-problemphyxmini0849-answerchoice}
  Resolution implied by choices spaced in hundreds of newtons per coulomb
\end{definition}
\begin{proof}
  This is the declared model definition of displayed Resolution In Newtons Per Coulomb.
\end{proof}
\begin{definition}[recorded Dataset Answer]
  \label{def:physics:phyx-mini-0849:phyxminiproblems-problemphyxmini0849-recordeddatasetanswer}
  \lean{PhyXMiniProblems.ProblemPhyXMini0849.recordedDatasetAnswer}
  \uses{def:physics:phyx-mini-0849:phyxminiproblems-problemphyxmini0849-answerchoice}
  The answer label recorded by the supplied dataset metadata
\end{definition}
\begin{proof}
  This is the declared model definition of recorded Dataset Answer.
\end{proof}
\begin{definition}[Rounds To Nearest Resolution]
  \label{def:physics:phyx-mini-0849:phyxminiproblems-problemphyxmini0849-roundstonearestresolution}
  \lean{PhyXMiniProblems.ProblemPhyXMini0849.RoundsToNearestResolution}
  A physical readout rounds to a display at a stated positive resolution
\end{definition}
\begin{proof}
  This is the declared model definition of Rounds To Nearest Resolution.
\end{proof}
\begin{definition}[Answer Matches Exterior Field]
  \label{def:physics:phyx-mini-0849:phyxminiproblems-problemphyxmini0849-answermatchesexteriorfield}
  \lean{PhyXMiniProblems.ProblemPhyXMini0849.AnswerMatchesExteriorField}
  \uses{def:physics:phyx-mini-0849:phyxminiproblems-problemphyxmini0849-conductingboxsetup, def:physics:phyx-mini-0849:phyxminiproblems-problemphyxmini0849-electricfieldstrengthinnewtonspercoulomb, def:physics:phyx-mini-0849:phyxminiproblems-problemphyxmini0849-answerchoice, def:physics:phyx-mini-0849:phyxminiproblems-problemphyxmini0849-displayedexteriorfieldstrengthinnewtonspercoulomb, def:physics:phyx-mini-0849:phyxminiproblems-problemphyxmini0849-displayedresolutioninnewtonspercoulomb, def:physics:phyx-mini-0849:phyxminiproblems-problemphyxmini0849-roundstonearestresolution}
  A choice matches the rounded strength immediately outside the box
\end{definition}
\begin{proof}
  This is the declared model definition of Answer Matches Exterior Field.
\end{proof}
\begin{definition}[Is Unique Matching Answer]
  \label{def:physics:phyx-mini-0849:phyxminiproblems-problemphyxmini0849-isuniquematchinganswer}
  \lean{PhyXMiniProblems.ProblemPhyXMini0849.IsUniqueMatchingAnswer}
  \uses{def:physics:phyx-mini-0849:phyxminiproblems-problemphyxmini0849-conductingboxsetup, def:physics:phyx-mini-0849:phyxminiproblems-problemphyxmini0849-answerchoice, def:physics:phyx-mini-0849:phyxminiproblems-problemphyxmini0849-answermatchesexteriorfield}
  A choice is the unique displayed match for the exterior field strength
\end{definition}
\begin{proof}
  This is the declared model definition of Is Unique Matching Answer.
\end{proof}
\begin{lemma}[point1 field strength rounding bounds]
  \label{lem:physics:phyx-mini-0849:phyxminiproblems-problemphyxmini0849-point1-field-strength-rounding-bounds}
  \lean{PhyXMiniProblems.ProblemPhyXMini0849.point1_field_strength_rounding_bounds}
  \uses{def:physics:phyx-mini-0849:phyxminiproblems-problemphyxmini0849-conductingboxsetup, def:physics:phyx-mini-0849:phyxminiproblems-problemphyxmini0849-matchesconductingboxscenario, def:physics:phyx-mini-0849:phyxminiproblems-problemphyxmini0849-matchesreportedsurfacedensitydata, def:physics:phyx-mini-0849:phyxminiproblems-problemphyxmini0849-matchessuppliedconductingboxfigure, def:physics:phyx-mini-0849:phyxminiproblems-problemphyxmini0849-usesstandardelectrostaticconstants, def:physics:phyx-mini-0849:phyxminiproblems-problemphyxmini0849-hasphysicalconductingboxparameters, def:physics:phyx-mini-0849:phyxminiproblems-problemphyxmini0849-satisfiesconductingboxelectrostatics, def:physics:phyx-mini-0849:phyxminiproblems-problemphyxmini0849-electricfieldstrengthinnewtonspercoulomb}
  The exterior field obtained from n e / epsilon\_0 lies within the rounding interval for 900 N/C. This is an intermediate derived result rather than an assumption
\end{lemma}
\begin{proof}
  The conclusion follows from the governing relations and figure readouts named in the statement of point1 field strength rounding bounds.
\end{proof}
% --- Archon declaration topology end ---
% --- Archon physics formalization source end ---
